% ============================================================
% Section 182: 两能级原子的Rabi振荡
% ============================================================
\section{量子力学：两能级原子的Rabi振荡}
\label{sec:rabi-oscillation}

\begin{modelbox}[题目：两能级原子的Rabi振荡]
两能级原子，能级差 $E_2-E_1=\hbar\omega_{21}$。初始处于基态 $|1\rangle$，受到电磁辐射 $\mathcal{E}(t)=\mathcal{E}_m(e^{i\omega t}+e^{-i\omega t})$ 作用。
\begin{enumerate}
    \item 若 $\omega=\omega_{21}$，计算一段时间后原子处于 $|2\rangle$ 的概率；
    \item 若 $\omega$ 只是近似等于 $\omega_{21}$，定性差别如何？重新计算上述概率。
\end{enumerate}
\end{modelbox}

\begin{tipbox}[考点快速分析]
\begin{itemize}
    \item \textbf{旋波近似}：保留 $e^{\pm i(\omega-\omega_{21})t}$，略去 $e^{\pm i(\omega+\omega_{21})t}$
    \item \textbf{Rabi频率}：$\Omega=|a_{21}|=|d_{21}\mathcal{E}_m|/\hbar$
    \item \textbf{共振}：完全振荡，$P_{1\to2}=\sin^2(\Omega t)$
    \item \textbf{失谐}：有效频率增大为 $\Lambda/2=\sqrt{\Delta^2+4\Omega^2}/2$，最大概率降低
\end{itemize}
\end{tipbox}

\begin{derivationbox}[标准解题过程]
\textbf{1. 基本方程}

设 $|\psi(t)\rangle=c_1(t)e^{-i\omega_1t}|1\rangle+c_2(t)e^{-i\omega_2t}|2\rangle$。微扰 $H'=-ex\mathcal{E}(t)$，矩阵元 $H'_{ij}=\hbar a_{ij}(e^{i\omega t}+e^{-i\omega t})$。

旋波近似下（$\omega\approx\omega_{21}$）：
\[
\dot{c}_1=-ia_{12}c_2e^{i(\omega-\omega_{21})t},\quad\dot{c}_2=-ia_{21}c_1e^{-i(\omega-\omega_{21})t}.
\]

记失谐 $\Delta=\omega_{21}-\omega$，$\Omega=|a_{21}|$。

\textbf{2. 共振情形 ($\omega=\omega_{21}$)}

$\Delta=0$，方程简化为 $\ddot{c}_2=-\Omega^2c_2$。由 $c_2(0)=0$，$\dot{c}_2(0)=-ia_{21}$：
\[
c_2(t)=-i\frac{a_{21}}{\Omega}\sin\Omega t.
\]

跃迁概率：
\begin{equation}
\boxed{P_{1\to2}(t)=|c_2(t)|^2=\sin^2(\Omega t).}
\end{equation}

原子以 \textbf{Rabi 频率} $\Omega$ 在基态与激发态之间完全振荡。

\textbf{3. 近共振情形 ($\omega\approx\omega_{21}$)}

消去 $c_1$ 得 $\ddot{c}_2+i\Delta\dot{c}_2+\Omega^2c_2=0$。特征根 $\lambda_\pm=\frac12(\Delta\pm\Lambda)$，$\Lambda=\sqrt{\Delta^2+4\Omega^2}$。

由初值条件：
\[
c_2(t)=-i\frac{2a_{21}}{\Lambda}e^{i\Delta t/2}\sin\frac{\Lambda t}{2}.
\]

跃迁概率：
\begin{equation}
\boxed{P_{1\to2}(t)=\frac{4\Omega^2}{\Delta^2+4\Omega^2}\sin^2\Bigl(\frac{\Lambda t}{2}\Bigr)=\frac{4\Omega^2}{\Delta^2+4\Omega^2}\sin^2\Bigl(\frac{\sqrt{\Delta^2+4\Omega^2}}{2}t\Bigr).}
\end{equation}

\textit{定性差别}：
\begin{itemize}
    \item 最大概率 $\dfrac{4\Omega^2}{\Delta^2+4\Omega^2}<1$，失谐越大越难完全跃迁；
    \item 振荡频率加快为 $\Lambda/2>\Omega$；
    \item $|\Delta|\gg\Omega$ 时概率极小（远离共振）。
\end{itemize}
\end{derivationbox}

\begin{formulabox}[答案汇总]
\begin{center}
\begin{tabular}{ll}
\toprule
\textbf{结论} & \textbf{结果} \\
\midrule
共振概率 & $P(t)=\sin^2(\Omega t)$ \\
近共振概率 & $P(t)=\dfrac{4\Omega^2}{\Delta^2+4\Omega^2}\sin^2\Bigl(\dfrac{\sqrt{\Delta^2+4\Omega^2}}{2}t\Bigr)$ \\
最大概率（失谐） & $<1$，随 $|\Delta|$ 增大而减小 \\
\bottomrule
\end{tabular}
\end{center}
\end{formulabox}

\begin{insightbox}[物理图像]
\begin{itemize}
    \item Rabi 振荡是量子光学和核磁共振的核心结果。频率 $\Omega$ 正比于耦合强度（电场振幅），反映了原子与光场的能量交换。
    \item 失谐时的 ``sin$^2$ 因子'' 是量子力学版本的拍频现象：原子与光场的相位差导致有效耦合减弱。
    \item 在超强场下，旋波近似失效，需用全量子理论（Jaynes-Cummings 模型）处理。
\end{itemize}
\end{insightbox}
